\subsection{关于诱导测度的积分}

设 X 是 T 的一个局部紧子空间，\(\mu\) 是 T 上的一个正测度，\(\mu_{X}\) 是 \(\mu\) 在 X 上诱导的测度（Ch. IV, §5, No.~7）。对每个 \(t \in T\) ，我们用下列方式定义 X 上的一个测度 \(\lambda_{t}\) ：\(\lambda_{t} = \varepsilon_{t}\) （若 \(t \in X\) ），\(\lambda_{t} = 0\) （若 \(t \in C X\) ）。对每个定义在 X 上的有限数值函数 g，有 \(\int g(x) d\lambda_{t}(x) = g(t)\) （若 \(t \in X\) ），而 \(\int g(x) d\lambda_{t}(x) = 0\) （若 \(t \in C X\) ）。于是，若 g 是 \(\mathcal{K}(X)\) 中的一个函数，则依 \(\mu_{X}\) 的定义，我们有

\[
\mu_ {\mathrm{X}} (g) = \int \langle g, \lambda_ {t} \rangle d \mu (t).\tag{1}
\]

这就是说，可以写

\[
\mu_ {\mathrm{X}} = \int \lambda_ {t} d \mu (t)\tag{2}
\]

（§3, No.~1）。

现在让我们定义一个由 T 到 X 中的映射 \(\pi\) ：置 \(\pi(t) = t\) （对 \(t \in \mathbf{X}\) ），并置 \(\pi(t) = t_0\) （对 \(t \in \mathbb{C} \mathbf{X}\) ），其中 \(t_0\) 是 \(\mathbf{X}\) 的任意一点；可以写 \(\lambda_t = \varphi_{\mathbf{X}}(t) \varepsilon_{\pi(t)}\) ，对每个 \(t \in \mathbf{T}\) 而言。映射 \(\pi\) 是 \(\mu\) -可测的，因为它在 \(\mathbf{X}\) 和 \(\mathbb{C} \mathbf{X}\) 上的限制都是可测的（Ch. IV, §5, No.~10, 命题 16）；由此立即得出，对 \((\pi, \varphi_{\mathbf{X}})\) 是 \(\mu\) -适应的（§4, No.~1）。于是我们有下列结果：

命题 1. --- 对每个定义在 X 上的数值函数 \(g \geqslant 0\) ，

\[
\int^ {\bullet} g d \mu_ {\mathrm{X}} = \int_ {\mathrm{X}} ^ {\bullet} g d \mu\tag{3}
\]

（关于记号 \(\int_{\mathbf{X}}^{\bullet}\) ，参见 §5, No.~3, 例）。

注意到前面的注记与 (2)，关系式 (3) 由 §4 的定理 1 得出。

推论 1. --- 对 X 的每个子集 B，\(\mu_{\mathrm{X}}^{\bullet}(\mathrm{B}) = \mu^{\bullet}(\mathrm{B})\) ；为使 B 是局部 \(\mu_{\mathrm{X}}\) -可忽略的，必须且只须 B 是局部 \(\mu\) -可忽略的。

推论 2. --- 设 M 是 T 的一个子集。若 \(\mu\) 集中于 M，则 \(\mu_{X}\) 集中于 \(M \cap X\) 。

推论 3. --- 为使测度 \(\mu_{X}\) 为零，必须且只须 \(X\) 是局部 \(\mu\) -可忽略的。

注. --- 若 S 是 \(\mu\) 的支集，则 \(S \cap X\) （它在 X 中是闭的）依推论 2 包含 \(\mu_{X}\) 的支集，但可以与它不同。例如，若 \(\mu\) 是一个弥散测度而 X 是缩成一点的子空间，则诱导测度 \(\mu_{X}\) 为零，因而其支集是空的。但要注意，若 X 在 T 中是开集，则 \(\mu_{X}\) 的支集等于 \(S \cap X\) 。

命题 2. --- 为使一个由 X 到拓扑空间中的映射 g 是 \(\mu_{X}\) -可测的，必须且只须 g 在 X 中是 \(\mu\) -可测的（\(\S5\) , No.~3, 例）。

这由 §4 的命题 3 得出。

推论. --- 为使 X 的一个子集 B 是 \(\mu_{X}\) -可测的，必须且只须 B 是 \(\mu\) -可测的。

定理 1. --- 设 g 是定义在 X 上、取值于 \(\overline{R}\) 或一个 Banach 空间中的函数。为使 g 是本质 \(\mu_{X}\) -可积的，必须且只须 \(\mathbf{g}\) 在 X 中是本质 \(\mu\) -可积的（§5, No.~3, 例），此时

\[
\int \mathbf {g} d \mu_ {\mathrm{X}} = \int_ {\mathrm{X}} \mathbf {g} d \mu .\tag{4}
\]

这由 §4 的定理 2 得出。

推论 1. --- 为使 X 的一个子集 B 是本质 \(\mu_{X}\) -可积的，必须且只须它是本质 \(\mu\) -可积的，此时 \(\mu_{\mathrm{X}}(\mathrm{B}) = \mu(\mathrm{B})\) 。

推论 2. --- 设 g 是定义在 T 上且局部 \(\mu\) -可积的复值函数；则 g 在 X 上的限制 \(g_{X}\) 是局部 \(\mu_{X}\) -可积的，且

\[
(g \cdot \mu) _ {\mathrm{X}} = g _ {\mathrm{X}} \cdot \mu_ {\mathrm{X}}.\tag{5}
\]

这由定理 1（应用于诸函数 \(fg_{\mathbf{X}}\) （\(f \in \mathcal{K}(\mathbf{X}; \mathbf{C})\) ））以及复测度在 \(\mathbf{X}\) 上诱导的测度的定义（Ch. IV, §5, No.~7）立即得出。

推论 3. --- 设 \(\theta\) 是 T 上的一个复测度；则

\[
| \theta | _ {\mathrm{X}} = | \theta_ {\mathrm{X}} |.\tag{6}
\]

置 \(|\theta| = \mu\) ，并取 \(g\) 为绝对值为 1 且使 \(\theta = g \cdot \mu\) 的复值函数（§5, No.~5, 定理 2 的推论 3）而应用推论 2；则 \(\theta_{\mathrm{X}} = g_{\mathrm{X}} \cdot \mu_{\mathrm{X}}\) ；但 \(g_{\mathrm{X}}\) 是绝对值为 1 的函数，于是公式 (6) 由 §5, No.~2 的命题 2 得出。

评注. --- a) 推论 3 已用另一种方法证明过（Ch. IV, §5, No.~7, 引理 3）。

\begin{enumerate}
\def\labelenumi{\alph{enumi})}
\setcounter{enumi}{1}
\tightlist
\item
  依据推论 3，命题 1 的推论 1、2、3，命题 2，定理 1 及其推论 1、2 都可立即推广到复测度。
\end{enumerate}

概说. --- 对每个函数 \(\mathbf{f}\) （相应地，\(\mathbf{g}\) ）（定义在 X （相应地，T ）上、取值于 Banach 空间 F 或 \(\overline{\mathbf{R}}\) 中），用 \(\zeta(\mathbf{f})\) （相应地，\(\rho(\mathbf{g})\) ）表示 \(\mathbf{f}\) 到 T 上的以 0 延拓（相应地，\(\mathbf{g}\) 在 X 上的限制）。于是 \(\zeta(\rho(\mathbf{g})) = \varphi_{\mathrm{X}} \cdot \mathbf{g}\) ，\(\rho(\zeta(\mathbf{f})) = \mathbf{f}\) 。用 \(\mu'\) 表示 T 上的测度 \(\varphi_{\mathrm{X}} \cdot \mu\) 。对每个 \(p \in [1, +\infty]\) ，命题 1 与命题 2 蕴含：\(\zeta\) 把 \(\overline{\mathcal{L}}_{\mathrm{F}}^{p}(\mathrm{X}, \mu_{\mathrm{X}})\) 映入 \(\overline{\mathcal{L}}_{\mathrm{F}}^{p}(\mathrm{T}, \mu')\) ，而 \(\rho\) 把 \(\overline{\mathcal{L}}_{\mathrm{F}}^{p}(\mathrm{T}, \mu')\) 映到 \(\overline{\mathcal{L}}_{\mathrm{F}}^{p}(\mathrm{X}, \mu_{\mathrm{X}})\) 上，两种情形都保持范数，且当 \(p = 1\) 时还保持积分（定理 1）；过渡到相联系的 Hausdorff 空间，我们得到两个互为逆的同构。类似地，若把 \(\zeta\) 和 \(\rho\) 应用于正的数值函数，则本质上积分保持不变（命题 1）。这样，若约定把 X 上的函数与在 X-T 上为零的 T 上的函数等同起来，并把测度 \(\mu_{X}\) 与测度 \(\mu'\) 等同起来，则关于诱导测度的问题就归结为关于由密度定义的测度的问题，后者已在 §5 中处理。这类论证也可应用于复测度，依据定理 1 的推论 3。
% semantic-fragments: legacy-input-seam
\relax{}
